
\section{光栅衍射}

\begin{definition}[光栅衍射]
光栅：在透光或反光性能上具有周期性空间结构的光学元件。它可以分成两类：
\begin{itemize}
\item 透射光栅：光从光栅中透过去。典型结构是：许多等宽的狭缝，彼此等距离地沿直线排列。
\item 反射光栅：光不透过去，而是在周期性结构表面反射出来。很多精密光谱仪用的是反射光栅。
\end{itemize}
\[
d=a+b
\]这个 \(d\) 叫做光栅常数。它表示相邻两条缝对应点之间的距离，也就是“缝距”。光栅常数通常在微米到几十微米量级。每一条狭缝自己会产生单缝衍射；许多狭缝之间又会发生多缝干涉。最后屏幕上的真实光强分布，是这两个效应共同决定的。条纹特点：
\begin{itemize}
\item 亮：主极大很强；
\item 细：主极大很窄；
\item 疏：相邻主极大之间隔得比较开，中间还有暗纹和次级明纹。
\end{itemize}
\end{definition}

\begin{figure}[htbp]
\centering
\includegraphics[width=0.6\textwidth]{pdf-exercises/13.png}
\caption{光栅衍射示意图}
\end{figure}

\begin{theorem}[多缝干涉效应：光栅方程]
先暂时把每条缝看成很细的线光源，只考虑不同缝发出的光到达同一点 \(P\) 时的相位差。相邻两缝对应点的光程差为：
\[\delta=(a+b)\sin\theta=d\sin\theta\]
当相邻两缝发出的光到达 \(P\) 点时相位相同，也就是光程差为整数个波长时，写出光栅方程（主极大公式）：
\[
d\sin\theta=\pm k\lambda,~~k=0,1,2,\cdots
\]
这里的 \(k\) 是光栅衍射的级次：\(k=0\)：中央主极大；\(k=1\)：第一级主极大；\(k=2\)：第二级主极大；正负号表示中央明纹两侧对称的条纹。当满足光栅方程时，各缝发出的光在 \(P\) 点干涉相长。如果有 \(N\) 条缝，每条缝单独在 \(P\) 点产生的振幅是 \(E_0\)，那么主极大方向上所有振动同相叠加，总振幅就是\(E_P=NE_0\)光强正比于振幅平方，因此
\[
I_P=N^2 I_0
\]所以主极大非常亮。
\end{theorem}

\begin{figure}[htbp]
\centering
\includegraphics[width=0.6\textwidth]{pdf-exercises/14.png}
\caption{单色光斜入射时的光栅方程}
\end{figure}
